\begin{landscape}
\thispagestyle{empty}

\section*{ANEXOS}
\addcontentsline{toc}{section}{ANEXOS}

\begin{center}
    %\textbf{ MATRIZ DE CONSISTENCIA DE LA INVESTIGACIÓN} %\\
    \small Título: Evaluación del efecto de un Framework de Ingeniería de Prompts basado en la norma ISO/IEC 25010 en la mantenibilidad y la deuda técnica del código generado por inteligencia artificial generativa
\end{center}

\noindent

\small
\setlength{\tabcolsep}{4pt}
\renewcommand{\arraystretch}{0.9}

\begin{adjustbox}{max width=\linewidth, max height=0.95\textheight, keepaspectratio}
\begin{tabularx}{\linewidth}{|Y|Y|Y|Y|Y|}
\hline
\rowcolor[gray]{0.9} 
\textbf{Problemas} & \textbf{Objetivos} & \textbf{Hipótesis} & \textbf{Variables} & \textbf{Metodología} \\ \hline

\textbf{Problema General:} \newline
¿Qué efecto tiene la implementación de un Framework de Ingeniería de Prompts basado en la norma ISO/IEC 25010 en la mantenibilidad y la deuda técnica del código generado por inteligencia artificial generativa? \newline

\textbf{Problemas Específicos:}
\begin{itemize}[leftmargin=*, nosep, itemsep=6pt]
    \item \textbf{PE1:} ¿Qué efecto tiene la aplicación de patrones arquitectónicos de prompts en la complejidad ciclomática de McCabe del código generado?
    \item \textbf{PE2:} ¿Qué efecto tiene la implementación de un pipeline de análisis estático en la reducción de la deuda técnica del software generado?
    \item \textbf{PE3:} ¿Qué efecto tiene la aplicación del protocolo de validación basado en ISO/IEC 25010 en el índice de mantenibilidad del código generado?
\end{itemize}
&
\textbf{Objetivo General:} \newline
Evaluar el efecto de la implementación de un Framework de Ingeniería de Prompts basado en la norma ISO/IEC 25010 en la mantenibilidad y la deuda técnica del código generado por inteligencia artificial generativa. \newline

\textbf{Objetivos Específicos:}
\begin{itemize}[leftmargin=*, nosep, itemsep=6pt]
    \item \textbf{OE1:} Evaluar el efecto de la aplicación de patrones arquitectónicos de prompts en la complejidad ciclomática de McCabe.
    \item \textbf{OE2:} Evaluar el efecto de la implementación de un pipeline de análisis estático en la reducción de la deuda técnica.
    \item \textbf{OE3:} Evaluar el efecto de la aplicación del protocolo de validación basado en ISO/IEC 25010 en el índice de mantenibilidad.
\end{itemize}
&
\textbf{Hipótesis General:} \newline
La implementación de un Framework de Ingeniería de Prompts basado en la norma ISO/IEC 25010 reduce significativamente la deuda técnica y mejora la mantenibilidad del código generado por inteligencia artificial generativa en comparación con métodos empíricos. \newline

\textbf{Hipótesis Específicas:}
\begin{itemize}[leftmargin=*, nosep, itemsep=6pt]
    \item \textbf{HIP1:} El código generado mediante el uso de patrones arquitectónicos de prompts presenta menor complejidad ciclomática que el código generado mediante métodos empíricos.
    \item \textbf{HIP2:} El código generado con un pipeline de análisis estático presenta menor deuda técnica que el generado sin dicho pipeline.
    \item \textbf{HIP3:} El código generado bajo el protocolo de validación basado en ISO/IEC 25010 presenta un mayor índice de mantenibilidad que el código generado sin dicho protocolo.
\end{itemize}
&
\textbf{Variable Independiente:}
\begin{itemize}[leftmargin=*, nosep]
    \item Aplicación del Framework de Ingeniería de Prompts
    \begin{itemize}[nosep]
        \item Nivel 1: Sin framework (método empírico)
        \item Nivel 2: Con framework estructurado
    \end{itemize}
\end{itemize}

\textbf{Variables Dependientes:}
\begin{itemize}[leftmargin=*, nosep]
    \item Mantenibilidad del Software:
    \begin{itemize}[nosep]
        \item Índice de mantenibilidad
        \item Complejidad ciclomática
    \end{itemize}
    \item Deuda Técnica:
    \begin{itemize}[nosep]
        \item Tiempo de remediación
        \item Densidad de defectos
    \end{itemize}
\end{itemize}
&
\textbf{Método:} Hipotético-deductivo. \newline
\textbf{Tipo:} Aplicada / Tecnológica. \newline
\textbf{Nivel:} Explicativo con enfoque experimental. \newline
\textbf{Diseño:} Cuasi-experimental (grupo control vs grupo experimental). \newline

\textbf{Población/Muestra:} \newline
Módulos de lógica para microservicios generados por modelos de lenguaje (n = 50 muestras por grupo). \newline

\textbf{Técnicas/Instrumentos:} \newline
Experimentación y análisis estático de código mediante SonarQube y Radon.\\ \hline

\end{tabularx}
\end{adjustbox}
\end{landscape}